% part: intuitionistic-logic
% chapter: propositions-as-types
% section: sequent-natural-deduction

\documentclass[../../../include/open-logic-section]{subfiles}

\begin{document}

\olfileid[id]{int}{pty}{snd}

\olsection{Deduksi Alami Sekuen}

Kita menulis $\Gamma \Sequent !A$ jika terdapat !!{derivation} deduksi alami
dengan $\Gamma$ sebagai asumsi yang !!{undischarged} dan $!A$ sebagai
kesimpulan; jika $\Gamma$ kosong, kita menulis $\Sequent !A$.

Kita menulis $\Gamma, !A_1, \dots, !A_n$ untuk
$\Gamma \cup \{!A_1, \dots, !A_n\}$, dan $\Gamma, \Delta$ untuk
$\Gamma \cup \Delta$.

Perhatikan bahwa jika kita mempunyai $\Gamma \Sequent !A \land !B$, yakni
!!a{derivation} dengan $\Gamma$ sebagai asumsi yang !!{undischarged} dan
$!A \land !B$ sebagai !!{formula} akhir, maka dengan menerapkan
\Elim{\land} di bagian bawah kita memperoleh !!a{derivation} dengan asumsi
yang !!{undischarged} sama dan $!A$ sebagai kesimpulan. Dengan kata lain,
jika $\Gamma \Sequent !A \land !B$, maka $\Gamma \Sequent !A$.
\begin{prooftree}
  \Axiom$\Gamma \fCenter !A \land !B$
  \RightLabel{$\Elim{\land}$}
  \UnaryInf$\Gamma \fCenter !A$
  \DisplayProof\qquad\bottomAlignProof
  \Axiom$\Gamma \fCenter !A \land !B$
  \RightLabel{$\Elim{\land}$}
  \UnaryInf$\Gamma \fCenter !B$
\end{prooftree}
Label $\Elim{\land}$ menunjukkan hubungannya dengan aturan bernama sama
dalam deduksi alami.

Demikian pula, andaikan kita mempunyai $\Gamma, !A \Sequent !B$, yakni
!!a{derivation} dengan asumsi yang !!{undischarged} $\Gamma, !A$ dan
!!{formula} akhir~$!B$. Jika kita menerapkan aturan \Intro{\lif}, kita
mempunyai !!a{derivation} dengan $\Gamma$ sebagai asumsi yang
!!{undischarged} dan $!A \lif !B$ sebagai !!{formula} akhir, yakni
$\Gamma \Sequent !A \lif !B$. Perhatikan bahwa dengan notasi ini, kita
!!{discharge} asumsi secara lebih eksplisit.
\begin{prooftree}
  \Axiom $\Gamma, !A \fCenter !B$
  \RightLabel{$\Intro{\lif}$}
  \UnaryInf $\Gamma \fCenter !A \lif !B$
\end{prooftree}

Dengan cara yang sama, kita dapat menarik kesimpulan dari aturan-aturan lain,
yang diperinci sebagai berikut:
\begin{gather*}
  \Axiom $\Gamma \fCenter !A$
  \Axiom $\Delta \fCenter !B$
  \RightLabel{$\Intro{\land}$}
  \BinaryInf $\Gamma,\Delta \fCenter !A \land !B$
  \DisplayProof\\
  \Axiom $\Gamma \fCenter !A \land !B$
  \RightLabel{$\Elim{\land}_1$}
  \UnaryInf $\Gamma \fCenter !A$
  \DisplayProof
  \qquad
  \Axiom $\Gamma \fCenter !A \land !B$
  \RightLabel{$\Elim{\land}_2$}
  \UnaryInf $\Gamma \fCenter !B$
  \DisplayProof
  \\
  \Axiom $\Gamma \fCenter !A$
  \RightLabel{$\Intro{\lor}_1$}
  \UnaryInf $\Gamma \fCenter !A \lor !B$
  \DisplayProof\qquad
  \Axiom $\Gamma \fCenter !B$
  \RightLabel{$\Intro{\lor}_2$}
  \UnaryInf $\Gamma \fCenter !A \lor !B$
  \DisplayProof\\
  \Axiom $\Gamma \fCenter !A \lor !B$
  \Axiom $\Delta, !A \fCenter !C$
  \Axiom $\Delta', !B \fCenter !C$
  \RightLabel{$\Elim{\lor}$}
  \TrinaryInf $\Gamma, \Delta, \Delta' \fCenter !C$
  \DisplayProof
  \\
  \Axiom $\Gamma, !A \fCenter !B$
  \RightLabel{$\Intro{\lif}$}
  \UnaryInf $\Gamma \fCenter !A \lif !B$
  \DisplayProof
  \qquad
  \Axiom $\Delta \fCenter !A \lif !B$
  \Axiom $\Gamma \fCenter !A$
  \RightLabel{$\Elim{\lif}$}
  \BinaryInf $\Gamma, \Delta \fCenter !B$
  \DisplayProof
  \\
  \Axiom $\Gamma \fCenter \lfalse$
  \RightLabel{$\FalseInt$}
  \UnaryInf $\Gamma \fCenter !A$
  \DisplayProof
\end{gather*}

Setiap asumsi dengan sendirinya merupakan !!a{derivation} bagi $!A$
dari~$!A$; dengan kata lain, kita selalu mempunyai $!A \Sequent !A$.

\begin{prooftree}
  \AxiomC{$ $}
  \UnaryInfC{$!A \fCenter !A$}
\end{prooftree}

Secara bersama-sama, aturan-aturan ini dapat dipandang sebagai kalkulus
tentang keberadaan !!{derivation}s deduksi alami. Aturan-aturan itu juga
dapat dipandang sebagai varian notasi deduksi alami, tempat setiap langkah
mencatat bukan hanya !!{formula} yang kita !!{derive}, melainkan juga asumsi
yang !!{undischarged} dan yang kita gunakan untuk !!{derive} formula tersebut.

Sebagai contoh, dengan $!B$ sebagai singkatan bagi
$(!A \lor (!A \lif \lfalse))\lif \lfalse$, kita memperoleh:
\begin{prooftree}
  \Axiom$!A \fCenter !A$
  \RightLabel{$\Intro{\lor}[1]$}
  \UnaryInf$!A \fCenter !A \lor (!A \lif \lfalse)$
  \Axiom$!B \fCenter !B$
  \RightLabel{$\Elim{\lif}$}
  \BinaryInf$!A, !B \fCenter \lfalse$
  \RightLabel{$\Intro{\lif}$}
  \UnaryInf$!B \fCenter !A \lif \lfalse$
  \RightLabel{$\Intro{\lor}[2]$}
  \UnaryInf$!B \fCenter !A \lor (!A \lif \lfalse)$
  \Axiom$!B \fCenter !B$
  \RightLabel{$\Elim{\lif}$}
  \BinaryInf$!B \fCenter \lfalse$
  \RightLabel{$\Intro{\lif}$}
  \UnaryInf$\fCenter !B \lif \lfalse$
\end{prooftree}

\end{document}
